\section{Dividing the work between people, code and agents}

Every system above makes a choice about which work goes to the language model, which to ordinary code, and which stays with a person.
Research on automation addressed this long before LLMs: it treats the degree of automation as a design choice, made separately for gathering information, analysing it, deciding and acting \citep{parasuraman2000model}, and warns that automating the routine parts of a job leaves people with the hardest parts and less practice at them \citep{bainbridge1983ironies}.

The most direct recent evidence comes from economics.
\citet{zhu2026human} ran 280 complete AI-assisted research projects on four datasets.
An unconstrained multi-agent pipeline produced a critical failure in 72\% of runs.
With the same model, agents and prompts, the failure rate fell to 16\% under three rules: the LLM reasons but does not do the data work, data preparation and estimation are done by deterministic code, and a person decides at three points, namely the research question, the identification strategy and publication.
The most common failures removed were hallucinated references, inconsistent interpretation and infeasible research questions, not coding errors.
The system tested is the authors' own and only one model was used, but the result matches what the engineering studies found independently: the language model should reason and propose, code should compute, and people should decide where decisions are hard to undo.
